\begin{filecontents*}[overwrite]{synthetic-rag-benchmark-pipeline.bib}
@inproceedings{kargupta2025claimspect,
  author    = {Kargupta, Priyanka and Tian, Runchu and Han, Jiawei},
  title     = {Beyond True or False: Retrieval-Augmented Hierarchical Analysis of Nuanced Claims},
  booktitle = {Proceedings of the 63rd Annual Meeting of the Association for Computational Linguistics},
  year      = {2025},
  pages     = {29664--29679},
  doi       = {10.18653/v1/2025.acl-long.1434},
  url       = {https://aclanthology.org/2025.acl-long.1434/}
}
@inproceedings{abbasiantaeb2026multiaspect,
  author    = {Abbasiantaeb, Zahra and Lupart, Simon and Aliannejadi, Mohammad},
  title     = {Generating Multi-Aspect Queries for Conversational Search},
  booktitle = {Proceedings of the 19th Conference of the European Chapter of the Association for Computational Linguistics},
  year      = {2026},
  pages     = {8201--8217},
  doi       = {10.18653/v1/2026.eacl-long.383},
  url       = {https://aclanthology.org/2026.eacl-long.383/}
}
@inproceedings{huang2025datagen,
  author    = {Huang, Yue and Wu, Siyuan and Gao, Chujie and Chen, Dongping and Zhang, Qihui and Wan, Yao and Zhou, Tianyi and Xiao, Chaowei and Gao, Jianfeng and Sun, Lichao and Zhang, Xiangliang},
  title     = {{DATAGEN}: Unified Synthetic Dataset via Large Language Models},
  booktitle = {International Conference on Learning Representations},
  year      = {2025}
}
@misc{driouich2025diverseprivate,
  author       = {Driouich, Ilias and Cao, Hongliu and Thomas, Eoin},
  title        = {Diverse And Private Synthetic Datasets Generation for {RAG} Evaluation: A Multi-Agent Framework},
  year         = {2025},
  eprint       = {2508.18929},
  archivePrefix= {arXiv},
  url          = {https://arxiv.org/abs/2508.18929}
}
@inproceedings{shen2026dragon,
  author    = {Shen, Haiyang and Yan, Hang and Xing, Zhongshi and Liu, Mugeng and Li, Yue and Chen, Zhiyang and Wang, Yuxiang and Wang, Jiuzheng and Ma, Yun},
  title     = {{DRAGON}: Domain-specific Robust Automatic Data Generation for {RAG} Optimization},
  booktitle = {Findings of the Association for Computational Linguistics: EACL},
  year      = {2026},
  pages     = {1065--1078},
  doi       = {10.18653/v1/2026.findings-eacl.56},
  url       = {https://aclanthology.org/2026.findings-eacl.56/}
}
@misc{lee2025mhts,
  author       = {Lee, Jeongsoo and Kwon, Daeyong and Jin, Kyohoon and Jeong, Junnyeong and Sim, Minwoo and Kim, Minwoo},
  title        = {{MHTS}: Multi-Hop Tree Structure Framework for Generating Difficulty-Controllable {QA} Datasets for {RAG} Evaluation},
  year         = {2025},
  eprint       = {2504.08756},
  archivePrefix= {arXiv},
  doi          = {10.48550/arXiv.2504.08756},
  url          = {https://arxiv.org/abs/2504.08756}
}
@inproceedings{lin2025scirgen,
  author    = {Lin, Junyong and Dai, Lu and Han, Ruiqian and Sui, Yijie and Wang, Ruilin and Sun, Xingliang and Wu, Qinglin and Feng, Min and Liu, Hao and Xiong, Hui},
  title     = {{ScIRGen}: Synthesize Realistic and Large-Scale {RAG} Dataset for Scientific Research},
  booktitle = {Proceedings of the 31st ACM SIGKDD Conference on Knowledge Discovery and Data Mining},
  year      = {2025},
  pages     = {5619--5630},
  doi       = {10.1145/3711896.3737432},
  url       = {https://dl.acm.org/doi/10.1145/3711896.3737432}
}
@inproceedings{rahmani2024synthetic,
  author    = {Rahmani, Hossein A. and Craswell, Nick and Yilmaz, Emine and Mitra, Bhaskar and Campos, Daniel},
  title     = {Synthetic Test Collections for Retrieval Evaluation},
  booktitle = {Proceedings of the 47th International ACM SIGIR Conference on Research and Development in Information Retrieval},
  year      = {2024},
  pages     = {2647--2651},
  doi       = {10.1145/3626772.3657942},
  url       = {https://dl.acm.org/doi/10.1145/3626772.3657942}
}
@inproceedings{kenneweg2024ragval,
  author    = {Kenneweg, Tristan and Kenneweg, Philip and Hammer, Barbara},
  title     = {{RAGVAL}: Automatic Dataset Creation and Evaluation for {RAG} Systems},
  booktitle = {2024 2nd International Conference on Foundation and Large Language Models},
  year      = {2024},
  pages     = {470--475},
  doi       = {10.1109/FLLM63129.2024.10852482}
}
@inproceedings{jadon2025domainspecific,
  author    = {Jadon, Aryan and Patil, Avinash and Kumar, Shashank},
  title     = {Enhancing Domain-Specific Retrieval-Augmented Generation: Synthetic Data Generation and Evaluation using Reasoning Models},
  booktitle = {2025 IEEE International Conference on Industry 4.0, Artificial Intelligence, and Communications Technology},
  year      = {2025},
  pages     = {445--452},
  doi       = {10.1109/IAICT65714.2025.11100564}
}
@inproceedings{kim2025agorabench,
  author    = {Kim, Seungone and Suk, Juyoung and Yue, Xiang and Viswanathan, Vijay and Lee, Seongyun and Wang, Yizhong and Gashteovski, Kiril and Lawrence, Carolin and Welleck, Sean and Neubig, Graham},
  title     = {Evaluating Language Models as Synthetic Data Generators},
  booktitle = {Proceedings of the 63rd Annual Meeting of the Association for Computational Linguistics},
  year      = {2025},
  pages     = {6385--6403},
  url       = {https://github.com/neulab/data-agora}
}
@inproceedings{long2024survey,
  author    = {Long, Lin and Wang, Rui and Xiao, Ruixuan and Zhao, Junbo and Ding, Xiao and Chen, Gang and Wang, Haobo},
  title     = {On {LLMs}-Driven Synthetic Data Generation, Curation, and Evaluation: A Survey},
  booktitle = {Findings of the Association for Computational Linguistics: ACL},
  year      = {2024},
  pages     = {11065--11082},
  doi       = {10.18653/v1/2024.findings-acl.658},
  url       = {https://aclanthology.org/2024.findings-acl.658/}
}
@article{pastrian2026evolutionary,
  author  = {Pastri{\'a}n, Diego and Hidalgo, Nicol{\'a}s and Reyes, V{\'i}ctor and Rosas, Erika},
  title   = {Evolutionary Multi-Objective Prompt Learning for Synthetic Text Data Generation with Black-Box Large Language Models},
  journal = {Applied Sciences},
  year    = {2026},
  volume  = {16},
  number  = {8},
  pages   = {3623},
  doi     = {10.3390/app16083623}
}
@misc{cosentino2025healthbranches,
  author       = {Cosentino, Cristian and Defilippo, Annamaria and Dossena, Marco and Irwin, Christopher and Joubbi, Sara and Li{\`o}, Pietro},
  title        = {HealthBranches: Synthesizing Clinically-Grounded Question Answering Datasets via Decision Pathways},
  year         = {2025},
  eprint       = {2508.07308},
  archivePrefix= {arXiv},
  doi          = {10.48550/arXiv.2508.07308}
}
@misc{gong2026meddialogrubrics,
  author       = {Gong, Lecheng and Fang, Weimin and Yang, Ting and Tao, Dongjie and Guo, Chunxiao and Wei, Peng and Xie, Bo and Guan, Jinqun and Chen, Zixiao and Shi, Fang and Gu, Jinjie and Liu, Junwei},
  title        = {MedDialogRubrics: A Comprehensive Benchmark and Evaluation Framework for Multi-turn Medical Consultations in Large Language Models},
  year         = {2026},
  eprint       = {2601.03023},
  archivePrefix= {arXiv},
  doi          = {10.48550/arXiv.2601.03023}
}
@inproceedings{yue2021cliniqg4qa,
  author    = {Yue, Xiang and Zhang, Xinliang Frederick and Yao, Ziyu and Lin, Simon and Sun, Huan},
  title     = {{CliniQG4QA}: Generating Diverse Questions for Domain Adaptation of Clinical Question Answering},
  booktitle = {2021 IEEE International Conference on Bioinformatics and Biomedicine},
  year      = {2021},
  pages     = {580--587},
  doi       = {10.1109/BIBM52615.2021.9669300}
}
@misc{mehandru2026erreason,
  author       = {Mehandru, Nikita and Golchini, Niloufar and Garg, Namrata and LeSaint, Kathy T. and Nash, Christopher J. and Ramachandran, Anu and Zack, Travis and McCoy, Liam G. and Rodman, Adam and Bamman, David and Molina, Melanie and Alaa, Ahmed},
  title        = {{ER-Reason}: A Benchmark Dataset for {LLM} Clinical Reasoning in the Emergency Room},
  year         = {2026},
  eprint       = {2505.22919},
  archivePrefix= {arXiv},
  doi          = {10.48550/arXiv.2505.22919}
}
@misc{ziletti2025cohorts,
  author       = {Ziletti, Angelo and D'Ambrosi, Leonardo},
  title        = {Generating Patient Cohorts from Electronic Health Records using Two-Step Retrieval-Augmented Text-to-{SQL} Generation},
  year         = {2025},
  eprint       = {2502.21107},
  archivePrefix= {arXiv},
  doi          = {10.48550/arXiv.2502.21107}
}
@article{thio2026medigraf,
  author  = {Thio, Samuel and Lewis, Matthew and Denaxas, Spiros and Dobson, Richard J. B.},
  title   = {Unlocking Electronic Health Records: A Hybrid Graph {RAG} Approach to Safe Clinical {AI} for Patient {QA}},
  journal = {Frontiers in Digital Health},
  year    = {2026},
  volume  = {8},
  pages   = {1780700},
  doi     = {10.3389/fdgth.2026.1780700}
}
@article{rao2026scoping,
  author  = {Rao, Hanshu and Liu, Weisi and Wang, Haohan and Huang, I-Chan and He, Zhe and Huang, Xiaolei},
  title   = {A Scoping Review of Synthetic Data Generation by Language Models in Biomedical Research and Application: Data Utility and Quality Perspectives},
  journal = {Journal of Healthcare Informatics Research},
  year    = {2026},
  volume  = {10},
  pages   = {367--392},
  doi     = {10.1007/s41666-026-00229-9}
}
@article{montenegro2025medicalsynthetic,
  author  = {Montenegro, Larissa and Gomes, Luis M. and Machado, Jos{\'e} M.},
  title   = {What We Know About the Role of Large Language Models for Medical Synthetic Dataset Generation},
  journal = {AI},
  year    = {2025},
  volume  = {6},
  number  = {6},
  pages   = {109},
  doi     = {10.3390/ai6060109}
}
@article{sufi2024datascarcity,
  author  = {Sufi, Fahim},
  title   = {Addressing Data Scarcity in the Medical Domain: A {GPT}-Based Approach for Synthetic Data Generation and Feature Extraction},
  journal = {Information},
  year    = {2024},
  volume  = {15},
  number  = {5},
  pages   = {264},
  doi     = {10.3390/info15050264}
}
@misc{nadas2025synthetic,
  author = {Nadas and others},
  title  = {Synthetic Data Generation Using Large Language Models: Advances in Text and Code},
  year   = {2025}
}
@article{figueira2022survey,
  author  = {Figueira, Alvaro and Vaz, Bruno},
  title   = {Survey on Synthetic Data Generation, Evaluation Methods and GANs},
  journal = {Mathematics},
  year    = {2022},
  volume  = {10},
  number  = {15},
  pages   = {2733},
  doi     = {10.3390/math10152733}
}
@inproceedings{li2023synthetic,
  author    = {Li, Zhuoyan and Zhu, Hangxiao and Lu, Zhuoran and Yin, Ming},
  title     = {Synthetic Data Generation with Large Language Models for Text Classification: Potential and Limitations},
  booktitle = {Proceedings of the 2023 Conference on Empirical Methods in Natural Language Processing},
  year      = {2023}
}
@misc{arzideh2026clinicalembeddings,
  author = {Arzideh and others},
  title  = {Improving Retrieval Augmented Generation for Health Care by Fine-Tuning Clinical Embedding Models},
  year   = {2026}
}
@misc{qureshi2026quality,
  author = {Qureshi and others},
  title  = {Improving Medical Data Quality via Synthetic Data Generation},
  year   = {2026}
}
@misc{villarreal2026healthsurveys,
  author = {Villarreal-Zegarra and Bellido-Boza},
  title  = {Generation of Synthetic Data in Health Surveys Using Large Language Models},
  year   = {2026}
}
@misc{zafar2024kimag,
  author = {Zafar and others},
  title  = {{KI-MAG}: A Knowledge-Infused Abstractive Question Answering System in Medical Domain},
  year   = {2024}
}
@misc{jin2022biomedicalqa,
  author = {Jin and others},
  title  = {Biomedical Question Answering: A Survey of Approaches and Challenges},
  year   = {2022}
}
\end{filecontents*}

\documentclass[11pt,a4paper]{article}
\usepackage[utf8]{inputenc}
\usepackage[T1]{fontenc}
\usepackage[margin=2.35cm]{geometry}
\usepackage{amsmath,amssymb}
\usepackage{array}
\usepackage{booktabs}
\usepackage{longtable}
\usepackage{enumitem}
\usepackage{xcolor}
\usepackage{natbib}
\usepackage[hidelinks]{hyperref}
\usepackage{parskip}
\usepackage{xurl}
\usepackage{microtype}
\usepackage{listings}
\usepackage{caption}

\sloppy
\definecolor{bettergrey}{gray}{0.95}
\definecolor{betterblue}{RGB}{33, 65, 196}
\definecolor{bettergreen}{RGB}{28, 130, 68}
\definecolor{betterred}{RGB}{165, 45, 45}
\newcommand{\field}[1]{\texttt{#1}}
\newcommand{\stage}[2]{\subsection{\textcolor{betterblue}{#1}: #2}}
\newcommand{\sourceidea}[2]{\textbf{#1}: #2}
\newcommand{\yes}{\textcolor{bettergreen}{yes}}
\newcommand{\no}{\textcolor{betterred}{no}}

\lstdefinestyle{promptstyle}{
  basicstyle=\ttfamily\small,
  breaklines=true,
  frame=single,
  backgroundcolor=\color{bettergrey},
  columns=fullflexible,
  keepspaces=true
}

\title{\textbf{Synthetic Dataset Construction Pipeline}\\
\large A Controlled Multi-Aspect Clinical RAG Benchmark}
\author{Master Thesis Design Note}
\date{\today}

\begin{document}
\maketitle

\begin{abstract}
	This document proposes a complete synthetic dataset generation pipeline for a retrieval-augmented generation (RAG) benchmark. The benchmark starts from a hidden structured plan: query facets and answer structure are defined first, atomic clinical facts are generated from that plan, facts are rendered into narrative retrieval chunks, redundant gold clusters and hard distractor clusters are created, and qrels plus gold answers are derived directly from the fact--chunk--facet mapping. MIMIC-IV is used only as a reference for terminology, condition families, comorbidity co-occurrences, demographic bins, and clinical-note style; it is not the evidence corpus being annotated after the fact. The benchmark should be presented as a controlled multi-aspect retrieval stress test in a clinical style, not as a claim that LLM-generated records are a faithful substitute for real hospital data.
\end{abstract}

\tableofcontents
\newpage

\section{Pipeline Core Design and Objective}

The experimental objective is to construct a corpus, query set, gold evidence set, and evaluation protocol in which the following hypothesis can be tested cleanly:

\begin{quote}
	For broad clinical information needs that require evidence from several semantically distinct but individually redundant subtopics, a facility-location coverage objective retrieves more complete evidence than pure similarity top-$k$ and dispersion-only MMR.
\end{quote}

The original MIMIC-IV discharge-note pipeline did not yield a strong result because the corpus was not generated to contain this structure. Discharge notes are clinically rich, but they are noisy, patient-specific, stylistically inconsistent, and often dominated by copied sections or local admission details. In that setting, the benchmark asks a coverage method to find structure that may not actually be present in a clean way.

The new pipeline reverses the old MIMIC-IV workflow. Instead of starting from raw notes and trying to infer a benchmark, it starts from structured metadata about the dataset to be generated:

\begin{enumerate}[leftmargin=*]
	\item define query facets and answer structure;
	\item generate atomic clinical facts that support each facet;
	\item render those facts into narrative clinical text chunks;
	\item generate redundant gold clusters and hard distractor clusters;
	\item generate natural multi-aspect queries from the same hidden plan;
	\item derive qrels and gold answers directly from the fact--chunk--facet mapping;
	\item validate the embedding geometry before final evaluation.
\end{enumerate}

This turns the dataset from a retrospective annotation problem into a prospective benchmark design problem.

\section{Literature Map: Which Papers Matter and Why}

The papers can be grouped by the role they should play in the new benchmark. Most generic RAG-generation papers start from an existing data source and construct QA pairs on top of it. This pipeline borrows their construction methods but changes the starting point: the real source of truth is the hidden schema that defines facts, facets, clusters, qrels, and gold answers.

\begin{longtable}{@{}p{2.25cm}p{3.65cm}p{4.2cm}p{4.6cm}@{}}
	\caption{Source-to-design mapping for the proposed pipeline.}\\
	\toprule
	\textbf{Source} & \textbf{Paper name} & \textbf{Useful idea} & \textbf{How it is used here}\\
	\midrule
	\endfirsthead
	\toprule
	\textbf{Source} & \textbf{Paper name} & \textbf{Useful idea} & \textbf{How it is used here}\\
	\midrule
	\endhead
	\bottomrule
	\endfoot
	\citet{kargupta2025claimspect} & \emph{Beyond True or False} & Nuanced claims should be decomposed into aspect and sub-aspect hierarchies; retrieval can be aspect-discriminative. & Build a clinical facet hierarchy for each query, e.g.\ primary disease $\rightarrow$ subgroup $\rightarrow$ treatment duration $\rightarrow$ outcome. Each gold facet is an explicit node.\\
	\citet{abbasiantaeb2026multiaspect} & \emph{Generating Multi-Aspect Queries for Conversational Search} & Complex user needs can be represented as several distinct aspect queries rather than one monolithic rewrite. & Generate both a natural query and hidden subqueries, one per facet. Use subqueries for diagnostics and optional query expansion, not for the main retrieval baseline unless explicitly tested.\\
	\citet{lee2025mhts} & \emph{MHTS} & Answer-first multi-hop construction; difficulty depends on both hop count and semantic spread of evidence. & Generate answer/evidence plans before queries. Control difficulty by number of facets, evidence clusters, and embedding distances among supporting chunks.\\
	\citet{huang2025datagen} & \emph{DATAGEN} & Attribute-guided generation, group checking, validation, and user-specified constraints improve diversity and controllability. & Use structured attributes to guide case generation. Use role-aware group checking: preserve planned redundancy inside clusters while removing accidental duplicates across facets.\\
	\citet{driouich2025diverseprivate} & \emph{Diverse And Private Synthetic Datasets Generation for RAG Evaluation} & Diversity agent clusters source material and samples representatives; QA curation agent produces diverse QA pairs. & Use a diversity controller that monitors cluster sizes, semantic redundancy, and coverage of the attribute grid. Privacy is less central if using fully synthetic records, but pseudonymization remains useful if MIMIC exemplars are included.\\
	\citet{shen2026dragon} & \emph{DRAGON} & A RAG benchmark should connect documents, queries, answers, clues, citations, hop counts, and answerability levels. & Store explicit mappings from query $\rightarrow$ facets $\rightarrow$ gold chunks $\rightarrow$ fact IDs. Generate sentence- or fact-level citations.\\
	\citet{lin2025scirgen} & \emph{ScIRGen} & Query generation can be controlled through cognitive taxonomies and answer validation filters. & Stratify questions by cognitive operation: comparison, causal explanation, contrast, management pattern, outcome synthesis.\\
	\citet{rahmani2024synthetic} and \citet{kenneweg2024ragval} & \emph{Synthetic Test Collections} and \emph{RAGVAL} & LLM-generated test collections can be useful, but benchmarks must control internal-knowledge leakage and evaluation reliability. & Add a no-context or no-retrieval baseline; reject queries answerable from general medical knowledge alone.\\
	\citet{nadas2025synthetic}, \citet{figueira2022survey}, and \citet{li2023synthetic} & Synthetic data generation surveys and classification study & Useful synthetic data requires task-specific quality criteria, conditional/few-shot grounding, objective labels, anti-collapse checks, and authenticity reporting. & Treat diversity as a query-local structural property, not a prompt adjective. Prefer objectively labelable facets and log non-copying, prompt, model, seed, and decoding metadata.\\
	\citet{jadon2025domainspecific} & \emph{Enhancing Domain-Specific RAG} & Domain RAG evaluation benefits from fine-grained evidence spans and discontinuous reference spans. & Keep fact IDs and span offsets inside chunks when possible, not only chunk-level qrels.\\
	\citet{cosentino2025healthbranches} & \emph{HealthBranches} & Clinical QA can be generated from explicit decision pathways with full reasoning paths. & Use medical decision pathways or guideline-like templates as a source of clinically plausible treatment/outcome structure.\\
	\citet{gong2026meddialogrubrics} & \emph{MedDialogRubrics} & Synthetic patient agents should be grounded in atomic medical facts; evaluation can use fine-grained rubrics. & Generate each synthetic narrative only from an atomic fact table. Validate that every narrative claim is supported by that table.\\
	\citet{yue2021cliniqg4qa} & \emph{CliniQG4QA} & Diverse clinical questions can be induced by predicting/forcing question phrases. & Use question-type starters such as ``How do'', ``What differences'', ``Which factors'', ``When should'', and ``Why might'' to diversify query forms.\\
	\citet{ziletti2025cohorts} & \emph{Generating Patient Cohorts from EHRs} & Complex clinical criteria should be decomposed into structured inclusion/exclusion logic and standardized concepts. & Represent queries as cohort-like logical plans over condition, demographic, comorbidity, intervention, and outcome axes.\\
	\citet{zafar2024kimag} & \emph{KI-MAG} & UMLS-style entities and triples can ground generated medical QA, even though the original system is not a retrieval-diversification benchmark. & Use concepts and relations as optional latent constraints for condition--comorbidity--treatment--outcome facts, then hide them from the evaluated retrievers.\\
	\citet{jin2022biomedicalqa} and \citet{arzideh2026clinicalembeddings} & Biomedical QA survey and clinical embedding fine-tuning & Biomedical QA often needs exact terminology, evidence grounding, and multi-document reasoning; clinical chunks can still contain many intertwined facts and require filtering. & Keep chunks narrow enough for objective facet labels, test both general and biomedical embeddings, and avoid treating embedding fine-tuning as the main thesis variable.\\
	\citet{qureshi2026quality} and \citet{villarreal2026healthsurveys} & Medical synthetic-data quality and persona/profile generation & Synthetic health data should be checked for accuracy, consistency, completeness, subgroup balance, protected-attribute artifacts, and distributional drift. & Build latent patient profiles only as generation scaffolds, split by profile/query family, and report subgroup/facet balance rather than claiming population realism.\\
	\citet{mehandru2026erreason} & \emph{ER-Reason} & Clinical reasoning can be framed as evidence accumulation across workflow stages. & Optional extension: longitudinal facets, where early presentation, treatment choice, and outcome are separate evidence clusters.\\
	\citet{thio2026medigraf} & \emph{Unlocking Electronic Health Records} & Hybrid structured and unstructured clinical retrieval is useful for patient journey QA. & Keep both structured metadata and narrative chunks; use metadata for generation and gold annotation, not for the main unrestricted retrieval evaluation.\\
	\citet{rao2026scoping} and \citet{montenegro2025medicalsynthetic} & \emph{Scoping Review} and \emph{Role of LLMs for Medical Synthetic Dataset Generation} & Biomedical synthetic data generation requires explicit quality, utility, and validation reporting. & Include a generation audit, plausibility checks, and a transparent threat-to-validity section.\\
	\citet{kim2025agorabench} & \emph{Evaluating Language Models as Synthetic Data Generators} & A model's ability to generate synthetic data is not identical to its ability to solve tasks; generation quality should be measured. & Compare at least two generator models or run quality checks rather than assuming the largest local model is the best generator.\\
	\citet{pastrian2026evolutionary} & \emph{Evolutionary Multi-Objective Prompt Learning} & Multi-objective prompt optimization can balance fidelity and diversity. & Optional later phase: automatically tune prompts against clinical plausibility, objective label consistency, and cluster diversity objectives.\\
\end{longtable}

\section{Why the New Benchmark Should Not Be ``Synthetic Discharge Notes Only''}

There are three candidate strategies:

\begin{enumerate}[leftmargin=*]
	\item \textbf{Use MIMIC-IV directly.} This is realistic but too noisy and weakly controlled. The old experiment already suggests that natural discharge chunks do not reliably produce the desired cluster structure.
	\item \textbf{Generate full synthetic EHR records.} This is clinically familiar, but unnecessary for the retrieval hypothesis. Full records add noise, irrelevant sections, and medical-realism obligations that distract from the core coverage experiment.
	\item \textbf{Generate structured synthetic clinical evidence cards with narrative renderings.} This is the recommended approach. Each chunk is grounded in a structured fact table, but written in a discharge-note or clinical-summary style. This gives enough medical texture for realistic embeddings while preserving exact gold labels.
\end{enumerate}

The third approach best fits the thesis. It lets the benchmark contain:

\begin{itemize}[leftmargin=*]
	\item controlled redundant clusters;
	\item explicit complementary facets;
	\item hard non-gold distractors;
	\item exact fact-level gold evidence;
	\item query difficulty labels;
	\item generation provenance.
\end{itemize}

\section{Benchmark Objects and Data Model}

The benchmark should produce six primary tables. These tables are more important than the natural-language text itself because they make the evaluation reproducible. A generation-audit table is useful as an implementation artifact, but it is not part of the minimal benchmark object model.

\begin{longtable}{p{4.9cm}p{10.1cm}}
	\caption{Recommended benchmark artifacts.}\\
	\toprule
	\textbf{Artifact} & \textbf{Purpose and core fields}\\
	\midrule
	\endfirsthead
	\toprule
	\textbf{Artifact} & \textbf{Purpose and core fields}\\
	\midrule
	\endhead
	\bottomrule
	\endfoot
	\field{clinical\_facts.parquet} & Hidden structured source of truth. Fields: \field{fact\_id}, \field{condition}, \field{age\_group}, \field{comorbidities}, \field{interventions}, \field{duration}, \field{outcomes}, \field{complications}, \field{must\_mention}, \field{must\_not\_mention}, and \field{facet\_ids}.\\
	\field{chunks.parquet} & Retrieval corpus. Fields: \field{chunk\_id}, \field{fact\_ids}, \field{facet\_ids}, \field{query\_id}, \field{cluster\_id}, \field{cluster\_role}, \field{distractor\_type}, \field{text}, \field{style}, \field{split}, and optional \field{span\_offsets}.\\
	\field{queries.parquet} & Natural benchmark questions. Fields: \field{query\_id}, \field{query\_text}, \field{condition}, \field{query\_type}, \field{difficulty}, \field{n\_facets}, and \field{split}.\\
	\field{gold\_answers.parquet} & Canonical answer object. Fields: \field{query\_id}, \field{answer\_text}, facet-level answer summaries, answer claims, supporting fact IDs, and supporting chunk IDs.\\
	\field{query\_plans.parquet} & Hidden plan for each query. Fields: selected facets, answer outline, split, difficulty, cluster roles, distractor families, and generation constraints.\\
	\field{qrels.parquet} & Gold evidence. Fields: query--chunk--facet labels, relevance grade, support type, fact IDs, and optional span offsets.\\
\end{longtable}

Optional implementation output: \field{generation\_audit.parquet}, containing prompt hashes, model versions, decoding parameters, seeds, validators, rejection reasons, non-copying checks, and schema versions.

\section{Pipeline Overview}

The full pipeline follows the same eleven stages as the reduced version:

\begin{enumerate}[leftmargin=*]
	\item Build a facet catalogue.
	\item Generate hidden query plans.
	\item Generate structured fact rows.
	\item Render fact rows into text chunks via LLM.
	\item Create role-aware gold clusters.
	\item Generate hard distractor chunks.
	\item Generate the natural query and gold answer via LLMs.
	\item Create qrels directly from the hidden mapping.
	\item Validate the generated objects.
	\item Embed and filter by retrieval geometry.
	\item Evaluate retrieval methods.
\end{enumerate}

The key methodological shift is that the generation pipeline owns the gold labels. Gold evidence should not be reconstructed after the fact by asking an LLM to read arbitrary chunks. The LLM may validate, refine, or reject data, but the authoritative label should come from the hidden structured plan.

\section{Detailed Pipeline}

\stage{Stage 1}{Build a Facet Catalogue}

\paragraph{Goal.}
Define the benchmark scope, split policy, and clinical facet catalogue before generating any text.

\paragraph{Recommended claim.}
The benchmark evaluates retrieval diversity methods under controlled multi-aspect clinical information needs. It is not a medical decision-support dataset and should not be used to infer clinical policy.

\paragraph{Design choices.}
Use a clinical style because it naturally supports multi-aspect queries:

\begin{itemize}[leftmargin=*]
	\item disease plus comorbidity;
	\item disease plus demographic subgroup;
	\item treatment plus complication;
	\item intervention plus duration;
	\item outcome plus discharge disposition;
	\item disease trajectory over time.
\end{itemize}

\paragraph{Recommended scale.}
For a first thesis-quality benchmark:

\begin{itemize}[leftmargin=*]
	\item 4--8 primary conditions;
	\item 8--15 attribute axes;
	\item 400--800 queries total;
	\item 20,000--80,000 chunks;
	\item 3--5 gold facets per query;
	\item 5--20 gold chunks per facet;
	\item 10--50 distractor chunks per query in the top-$N$ candidate pool.
\end{itemize}

\paragraph{Splits.}
Use two reported splits:

\begin{itemize}[leftmargin=*]
	\item \field{validation}: used after the generation pipeline is fixed; choose retrieval hyperparameters, especially MMR and facility-location $\lambda$ values;
	\item \field{test}: final locked evaluation.
\end{itemize}

Do not tune filtering thresholds or retrieval hyperparameters on the test split. Split at the \field{query\_family} or latent-profile level, not only at the row level: surface variants of the same query plan, chunks generated from the same patient profile, and near-duplicate distractor templates should stay in the same split. This prevents leakage where validation and test differ only by wording.

\paragraph{MIMIC-IV and optional clinical scaffolds.}

MIMIC-IV can be used as a reference for terminology, frequent conditions, comorbidity distributions, age groups, and note style. It should not be treated as the corpus from which gold labels are discovered.

There are four viable ways to seed the catalogue.

\paragraph{Option A: MIMIC-seeded synthetic benchmark.}
Use MIMIC-IV only to estimate plausible distributions and wording:

\begin{itemize}[leftmargin=*]
	\item frequent conditions;
	\item common comorbidities;
	\item age distributions;
	\item common treatments;
	\item discharge-note style;
	\item typical section headers.
\end{itemize}

Then generate a new synthetic corpus from the structured schema. This is the recommended default because it gives realism without surrendering control.

\paragraph{Option B: Guideline or decision-pathway seeded benchmark.}
Use decision pathways, guideline-like rules, or medical educational material to define treatment and outcome structures. This follows the HealthBranches logic of grounding questions in explicit decision paths \citep{cosentino2025healthbranches}. It is useful when you want the clinical facts to be medically coherent without using patient-level EHR data.

\paragraph{Option C: Fully synthetic from scratch.}
Define a hand-built ontology and ask the LLM to populate it. This maximizes control but requires stronger validation because clinical plausibility depends heavily on the generator.

\paragraph{Option D: Concept- or KG-seeded benchmark.}
Use UMLS-style concepts, relation triples, or normalized clinical vocabularies as latent constraints for facts. This borrows the useful part of KI-MAG \citep{zafar2024kimag}: structured medical knowledge can constrain terminology and relations. The graph or concept layer should be used for generation and auditing only, not as an evaluation-time retrieval shortcut.

\paragraph{Recommendation.}
Use Option A plus a small amount of Option B and, where available, light Option D concept grounding. MIMIC supplies style and distributions; decision pathways supply management logic; concept grounding reduces terminology drift. The final corpus is still generated from the hidden schema.

\paragraph{Facet catalogue.}

\paragraph{Purpose.}
The facet ontology defines which aspects can appear in queries and which clusters should exist in the corpus. This is inspired by ClaimSpect's hierarchical aspect decomposition \citep{kargupta2025claimspect} and by cohort-style criteria decomposition \citep{ziletti2025cohorts}.

\paragraph{Example ontology.}

\begin{longtable}{p{3.5cm}p{11.5cm}}
	\caption{Example axes in the clinical facet ontology.}\\
	\toprule
	\textbf{Axis} & \textbf{Values}\\
	\midrule
	\endfirsthead
	\toprule
	\textbf{Axis} & \textbf{Values}\\
	\midrule
	\endhead
	\bottomrule
	\endfoot
	Primary condition & encephalitis/myelitis, heart failure, pneumonia, type 2 diabetes, chronic kidney disease, sepsis, stroke, COPD, liver cirrhosis, pulmonary embolism\\
	Demographics & age $<40$, age 40--64, age 65--75, age $>75$, male, female, race/ethnicity group if ethically justified and handled carefully\\
	Comorbidities & uncomplicated diabetes, complicated diabetes, renal disease, congestive heart failure, chronic pulmonary disease, dementia, malignancy, liver disease, immunosuppression\\
	Interventions & antibiotics, antivirals, steroids, anticoagulation, dialysis, oxygen escalation, rehabilitation, surgery, medication adjustment\\
	Durations & short course, prolonged course, delayed initiation, early transition, recurrent treatment\\
	Outcomes & home discharge, skilled nursing facility, inpatient rehab, readmission, functional improvement, persistent deficit, complication\\
	Reasoning relation & comparison, contrast, causal relation, temporal trigger, subgroup difference, management trade-off\\
\end{longtable}

\paragraph{Facet template.}
Each facet should be represented as a structured object:

\begin{lstlisting}[style=promptstyle]
{
  "facet_id": "F_encephalitis_elderly_rehab_outcome",
  "primary_condition": "encephalitis",
  "subgroup": {"age": ">75"},
  "clinical_axis": "rehabilitative_outcome",
  "answer_role": "comparison_arm",
  "minimum_gold_chunks": 6,
  "target_cluster_size": 15
}
\end{lstlisting}

\paragraph{Important constraint.}
Facets should be individually answerable but collectively necessary. A single chunk must not answer the whole query. This is what creates the multi-source retrieval problem.

Facets should also be objectively labelable from the hidden schema. Prefer labels such as age group, comorbidity presence, treatment-duration category, discharge disposition, complication type, rehabilitation setting, or explicitly stated response to therapy. Avoid facets that require ambiguous clinical judgment unless they are backed by an explicit rubric, because subjective labels make retrieval evaluation hard to interpret \citep{li2023synthetic}.

\stage{Stage 2}{Hidden Query Plans}

\paragraph{Purpose.}
Generate query plans from the catalogue before writing query text or evidence chunks. A query plan is the hidden recipe that says which facets must be covered, which clusters are gold, which cluster is intentionally dominant, which distractor families should be generated, and what the answer must contain. This follows the answer-first idea in MHTS \citep{lee2025mhts} and the query--answer--clue mapping in DRAGON \citep{shen2026dragon}.

\paragraph{Blueprint structure.}

\begin{lstlisting}[style=promptstyle]
{
  "query_id": "Q_000142",
  "primary_condition": "encephalitis_or_myelitis",
  "query_type": "subgroup_comparison",
  "facets": [
    "F_diabetes_treatment_duration",
    "F_diabetes_rehab_outcome",
    "F_elderly_treatment_duration",
    "F_elderly_rehab_outcome"
  ],
  "contrast_relation": "compare diabetes vs age_over_75",
  "expected_answer_outline": [
    "describe treatment duration in uncomplicated diabetes subgroup",
    "describe rehabilitative outcomes in uncomplicated diabetes subgroup",
    "describe treatment duration in age_over_75 subgroup",
    "describe rehabilitative outcomes in age_over_75 subgroup",
    "synthesize differences and caveats"
  ],
  "difficulty": {
    "n_facets": 4,
    "dominant_gold_cluster": "C1",
    "dominant_gold_chunks": 18,
    "complementary_gold_chunks_per_cluster": 8,
    "hard_distractor_chunks": 24,
    "expected_topk_failure": "oversample_dominant_gold_cluster"
  },
  "split": "validation"
}
\end{lstlisting}

The output of this stage is \field{query\_plans.parquet}. The important point is that diversity is already specified here: the plan says which clusters are gold, which one is intentionally redundant, which facets are complementary, and which distractor families should be generated.

\paragraph{Difficulty control.}
Assign a difficulty score:

\[
	D(q) = \alpha \cdot n_{\text{facets}} + \beta \cdot n_{\text{clusters}} + \gamma \cdot \overline{d}_{\text{facet}} + \delta \cdot n_{\text{distractor types}},
\]

where $\overline{d}_{\text{facet}}$ is the average semantic distance among gold facet clusters after embedding. This adapts MHTS's point that difficulty is not just hop count; it also depends on the semantic spread of evidence \citep{lee2025mhts}.

\stage{Stage 3}{Generate Structured Fact Rows}

\paragraph{Purpose.}
Create the hidden structured truth from which chunks and answers are generated. Each fact row says what the generated text is allowed to mention and which facet it supports. This stage follows MedDialogRubrics' principle of grounding patient simulation in atomic medical facts \citep{gong2026meddialogrubrics}.

\paragraph{Atomic fact example.}

\begin{lstlisting}[style=promptstyle]
{
  "fact_id": "FACT_003819",
  "patient_profile_id": "P_10482",
  "encounter_id": "E_22194",
  "primary_condition": "encephalitis",
  "concept_ids": ["C0014038", "C0009450"],
  "age": 79,
  "age_group": ">75",
  "comorbidities": ["hypertension", "atrial_fibrillation"],
  "intervention": "acyclovir",
  "treatment_duration": "14 days",
  "rehab_disposition": "inpatient rehabilitation",
  "outcome": "persistent gait instability but improved orientation",
  "facet_ids": [
    "F_elderly_treatment_duration",
    "F_elderly_rehab_outcome"
  ],
  "must_mention": [
    "14-day antiviral course",
    "age over 75",
    "inpatient rehabilitation",
    "gait instability"
  ],
  "must_not_mention": [
    "diabetes",
    "dialysis",
    "complete recovery"
  ]
}
\end{lstlisting}

For each facet, generate enough rows to create a cluster. These rows should vary across synthetic patients while preserving the same facet meaning. A structured fact row may contain fields useful for more than one answer claim, but rendered chunks should usually be scoped to one facet unless the benchmark intentionally includes an easy ablation. The output of this stage is \field{clinical\_facts.parquet}.

\paragraph{Generation method.}
Use attribute-guided generation \citep{huang2025datagen}. The prompt should specify:

\begin{itemize}[leftmargin=*]
	\item the primary condition;
	\item normalized concept IDs or relation triples when available;
	\item the demographic subgroup;
	\item comorbidities;
	\item allowed interventions;
	\item target outcome;
	\item the exact facet IDs the fact should support;
	\item forbidden facts to prevent leakage across facets.
\end{itemize}

\paragraph{Validation.}
Reject a generated fact if:

\begin{itemize}[leftmargin=*]
	\item required fields are missing;
	\item the clinical relation is internally inconsistent;
	\item it mentions a forbidden subgroup;
	\item it supports more facets than intended;
	\item it contains an impossible treatment/outcome pair.
\end{itemize}

\stage{Stage 4}{Render Fact Rows into Text Chunks via LLM}

\paragraph{Purpose.}
Render structured facts into realistic text chunks that can be embedded and retrieved.

\paragraph{Chunk styles.}
Generate multiple styles so the corpus is not one-note:

\begin{itemize}[leftmargin=*]
	\item discharge-summary brief hospital course;
	\item progress-note assessment and plan;
	\item rehab consultation summary;
	\item problem-list paragraph;
	\item outcome follow-up note;
	\item medication-management note.
\end{itemize}

\paragraph{Narrative generation prompt.}

\begin{lstlisting}[style=promptstyle]
You are generating a synthetic clinical evidence chunk for a RAG benchmark.
Use only the structured facts provided below.
Do not add diagnoses, treatments, durations, outcomes, or demographics that are
not present in the facts.

Style: brief hospital course problem-list paragraph.
Length: 90-160 words.
Required facet IDs: {facet_ids}
Required facts: {atomic_facts}
Forbidden facts: {forbidden_facts}

Write one self-contained clinical paragraph. It should sound like a deidentified
clinical note, but it must not include names, dates, addresses, or real identifiers.
\end{lstlisting}

\paragraph{Traceability requirement.}
Every chunk must retain:

\begin{itemize}[leftmargin=*]
	\item \field{fact\_ids}: which facts were used;
	\item \field{facet\_ids}: which facets the chunk supports;
	\item \field{query\_id}, \field{cluster\_id}, \field{cluster\_role}, and \field{split};
	\item \field{span\_offsets}: optional character offsets for fact mentions;
	\item \field{generation\_prompt\_id}: which prompt generated it.
\end{itemize}

The output of this stage is \field{chunks.parquet}.

\paragraph{Why not rely on LLM annotation afterward?}
Because the old pipeline used LLM map-reduce over candidate chunks to discover gold facets after retrieval. That makes gold annotation dependent on noisy text and candidate-pool choices. Here, the gold labels are known before embedding. LLMs can audit the labels, but they should not be the only source of truth.

\stage{Stage 5}{Create Role-Aware Gold Clusters}

\paragraph{Purpose.}
Facility-location wins when there are dense redundant regions and multiple complementary regions. Therefore, each facet must have a cluster of several relevant chunks.

\paragraph{Cluster design.}
For each query facet:

\begin{itemize}[leftmargin=*]
	\item generate 5--20 chunks that support that facet;
	\item vary patients, wording, and note style;
	\item keep the same facet-level answer;
	\item avoid generating exact paraphrases only;
	\item preserve planned within-cluster redundancy while removing accidental cross-facet duplication;
	\item ensure chunks within a facet are closer to each other than to chunks in other facets.
\end{itemize}

\paragraph{Role-aware deduplication.}
Duplicate filtering must not erase the experimental condition. Exact copies and accidental near-duplicates across facets should be rejected, but planned redundancy inside a dominant or complementary cluster should be preserved up to its target size. Otherwise the pipeline removes the very top-$k$ failure mode it is trying to test.

\paragraph{Example.}
For the query:

\begin{quote}
	For patients diagnosed with encephalitis or myelitis, how do treatment durations and rehabilitative outcomes differ between those with uncomplicated diabetes and elderly patients over age 75?
\end{quote}

generate these gold clusters:

\begin{enumerate}[leftmargin=*]
	\item encephalitis + uncomplicated diabetes + treatment duration;
	\item encephalitis + uncomplicated diabetes + rehab outcome;
	\item encephalitis + age $>75$ + treatment duration;
	\item encephalitis + age $>75$ + rehab outcome.
\end{enumerate}

Top-$k$ should be tempted to over-sample the most query-similar cluster, for example encephalitis + diabetes + treatment, because it shares many exact query terms. Facility-location should spend fewer picks in that redundant cluster and allocate representatives to the other clusters.

\stage{Stage 6}{Generate Hard Distractor Chunks}

\paragraph{Purpose.}
The benchmark must distinguish coverage from mere dispersion. MMR should not automatically win by selecting semantically distant chunks. Therefore, add distractors that are diverse but not relevant.

\paragraph{Distractor types.}

\begin{longtable}{p{3.5cm}p{6cm}p{5.2cm}}
	\caption{Distractor categories and expected retrieval behavior.}\\
	\toprule
	\textbf{Distractor type} & \textbf{Example} & \textbf{Expected effect}\\
	\midrule
	\endfirsthead
	\toprule
	\textbf{Distractor type} & \textbf{Example} & \textbf{Expected effect}\\
	\midrule
	\endhead
	\bottomrule
	\endfoot
	Same condition, wrong subgroup & encephalitis in middle-aged patients without diabetes & Top-$k$ may retrieve these because condition terms match.\\
	Same subgroup, wrong condition & elderly patients with pneumonia rehab outcomes & MMR may retrieve these because they add demographic/outcome diversity.\\
	Same intervention, wrong disease & acyclovir duration for shingles rather than encephalitis & Similar treatment vocabulary creates hard negatives.\\
	Same outcome, wrong clinical axis & rehab outcome after stroke, not encephalitis & Outcome terms match but facet is wrong.\\
	Broad textbook-style background & general encephalitis treatment principles & Tests whether the query is answerable only from corpus-specific evidence.\\
	Rare outlier & unusual complication unrelated to the query & MMR may overvalue outliers; facility-location should ignore isolated points.\\
\end{longtable}

\paragraph{Important distinction.}
The distractors should be semantically plausible, not nonsense. A trivial distractor does not test retrieval quality.

\stage{Stage 7}{Generate the Natural Query and Gold Answer via LLMs}

\paragraph{Purpose.}
Generate realistic broad clinical queries while preserving the hidden facet decomposition.

This stage draws from MQ4CS/MASQ's multi-aspect query framing \citep{abbasiantaeb2026multiaspect}, CliniQG4QA's question diversity idea \citep{yue2021cliniqg4qa}, and ScIRGen's use of cognitive/question taxonomies \citep{lin2025scirgen}.

\paragraph{Query types.}

\begin{itemize}[leftmargin=*]
	\item subgroup comparison: ``How do outcomes differ between A and B?''
	\item management trade-off: ``How is treatment adjusted when condition X co-occurs with Y?''
	\item temporal trigger: ``Which findings tend to trigger escalation from treatment A to B?''
	\item outcome synthesis: ``What discharge and rehabilitation patterns are reported across subgroups?''
	\item complication contrast: ``Which complications change the management pathway?''
\end{itemize}

\paragraph{Query generation prompt.}

\begin{lstlisting}[style=promptstyle]
Generate one broad clinical research-style question for a RAG benchmark.
The question must require evidence from all listed facets.
It must not reveal the facet IDs.
It must not be answerable from general medical knowledge alone; it should ask
about patterns in the provided synthetic corpus.

Primary condition: {condition}
Required facets:
{facet_descriptions}
Question type: {query_type}
Target difficulty: {difficulty}

Return JSON:
{
  "query_text": "...",
  "subqueries": ["...", "...", "..."],
  "why_multi_aspect": "...",
  "forbidden_single_source_answer": true
}
\end{lstlisting}

\paragraph{Subqueries.}
Store subqueries but do not use them for the main top-$k$/MMR/facility-location comparison unless running a separate multi-query retrieval ablation. The main experiment should rank from the single natural query so that methods face the same information need.

\paragraph{Gold answer construction.}
The answer is also produced at this stage, during benchmark construction, before embedding and before any retrieval method is evaluated. The source of truth is the hidden scaffold: \field{query\_plans.parquet}, \field{clinical\_facts.parquet}, and the fact-to-chunk mapping stored in \field{chunks.parquet}. This is close to DRAGON's explicit query--answer--clue--source mapping \citep{shen2026dragon}.

Use two separate prompts or procedures:

\begin{enumerate}[leftmargin=*]
	\item \textbf{Query prompt}: given only the hidden plan, write the user-facing \field{query\_text}; do not expose the answer wording.
	\item \textbf{Answer prompt}: given the required facets, the answer outline, and the structured fact rows grouped by facet, write a canonical answer. The prompt may receive \field{fact\_ids} and their derived \field{chunk\_ids} for citation, but it should not receive distractor chunks or retrieval rankings.
\end{enumerate}

\paragraph{Answer format.}

\begin{lstlisting}[style=promptstyle]
{
  "query_id": "Q_000142",
  "answer_text": "Across the synthetic corpus, patients with ...",
  "facet_answer_summaries": [
    {
      "facet_id": "F_diabetes_treatment_duration",
      "summary": "...",
      "required_chunk_ids": ["C_001", "C_017"]
    }
  ],
  "minimum_sufficient_evidence": {
    "one_chunk_per_facet": true,
    "required_facets": ["F1", "F2", "F3", "F4"]
  }
}
\end{lstlisting}

\paragraph{Avoiding leakage.}
The answer should not introduce new facts not present in \field{clinical\_facts.parquet}. Use an LLM judge only to check consistency:

\begin{itemize}[leftmargin=*]
	\item every answer claim is supported by at least one fact;
	\item every required facet is mentioned;
	\item no unsupported clinical recommendation is added;
	\item no distractor fact appears in the answer.
\end{itemize}

The output of this stage is \field{queries.parquet} and \field{gold\_answers.parquet}.

\stage{Stage 8}{Create Qrels Directly from the Hidden Mapping}

\paragraph{Purpose.}
Create \field{qrels.parquet} from the known mapping between queries, chunks, facts, and facets. The LLM-generated text is not re-annotated after retrieval to discover gold labels; qrels are derived from the hidden construction graph.

\paragraph{Example qrels.}

\begin{lstlisting}[style=promptstyle]
query_id  chunk_id  facet_id  relevance_grade  fact_ids
Q_001     CH_017    F3        1                [FACT_001]
Q_001     CH_017    F4        1                [FACT_001]
Q_001     CH_088    none      0                []
\end{lstlisting}

If a chunk is allowed to support multiple facets, every supported query--chunk--facet relationship must be represented explicitly. Distractor chunks receive no positive gold qrel for the target query, even when they are semantically close.

\stage{Stage 9}{Validate the Generated Objects}

\paragraph{Validation agents.}
Use a multi-agent curation pattern similar to DATAGEN and Driouich et al.\ \citep{huang2025datagen,driouich2025diverseprivate}:

\begin{enumerate}[leftmargin=*]
	\item \textbf{Schema validator}: checks JSON fields and allowed values.
	\item \textbf{Clinical plausibility validator}: checks that the case is medically plausible.
	\item \textbf{Facet-support validator}: checks whether the chunk supports exactly the intended facet(s).
	\item \textbf{Diversity validator}: computes similarity, checks cluster quality, and applies role-aware duplicate handling.
	\item \textbf{No-leakage validator}: checks that distractors do not accidentally answer a gold facet.
	\item \textbf{Non-copying validator}: checks that MIMIC-seeded style examples or guideline snippets were not copied into the synthetic corpus.
	\item \textbf{Re-extraction validator}: extracts schema fields back from generated chunks and compares them with the hidden facts.
	\item \textbf{Profile and subgroup validator}: checks facet balance, profile-level split integrity, and obvious protected-attribute artifacts.
	\item \textbf{No-context baseline validator}: checks whether a model can answer the query without retrieval.
\end{enumerate}

\paragraph{Automatic checks.}
Run these checks before embedding:

\begin{longtable}{p{4cm}p{6.1cm}p{4.7cm}}
	\caption{Recommended validation checks.}\\
	\toprule
	\textbf{Check} & \textbf{Metric} & \textbf{Default action}\\
	\midrule
	\endfirsthead
	\toprule
	\textbf{Check} & \textbf{Metric} & \textbf{Default action}\\
	\midrule
	\endhead
	\bottomrule
	\endfoot
	Structured completeness & all required fields non-null & reject object\\
	Query facet count & every query has 3--5 required facets & regenerate query plan\\
	Facet balance & gold chunks per facet in target range & regenerate underrepresented facets\\
	Cluster role consistency & every \field{cluster\_id} has exactly one \field{cluster\_role} & reject or repair cluster metadata\\
	Traceability & every chunk traces back to fact IDs and facet IDs & reject chunk\\
	Forbidden fact leakage & generated chunk mentions a \field{must\_not\_mention} fact & rewrite or reject chunk\\
	Near-duplicate chunks & cosine similarity above threshold, e.g.\ $>0.94$ & role-aware handling: preserve planned redundancy, reject exact copies and accidental cross-facet duplicates\\
	Authenticity/non-copying & overlap against any MIMIC seed or guideline source above threshold & rewrite or reject object\\
	Cluster quality & intra-facet similarity, inter-facet separation, silhouette or related score & regenerate weak clusters\\
	Distractor leakage & distractor judged relevant to required facet & relabel or reject\\
	Query answerability & every facet has at least one gold chunk & reject query\\
	Single-chunk shortcut & one chunk supports all required facets & reject or downgrade difficulty\\
	Canonical answer coverage & answer uses every required facet at least once & rewrite answer or reject query\\
	No-retrieval answerability & LLM without context answers too well & reject query or make it more corpus-specific\\
\end{longtable}

\paragraph{Why the no-context check matters.}
RAGVAL emphasizes that public or generic knowledge can make a RAG benchmark misleading because the generator may answer without retrieval \citep{kenneweg2024ragval}. In this benchmark, queries should ask about corpus-specific synthetic patterns, not universal facts such as ``What is encephalitis?''.

\stage{Stage 10}{Embed and Filter by Retrieval Geometry}

\paragraph{Purpose.}
The dataset should not be accepted merely because the hidden labels look good. After embedding, the candidate pools must exhibit the intended geometry. These are benchmark-construction checks, so they may use hidden fields such as \field{query\_id}, \field{facet\_ids}, \field{cluster\_id}, \field{cluster\_role}, \field{distractor\_type}, and \field{qrels.parquet}. These fields validate the dataset; they are not exposed to retrieval methods.

\paragraph{Procedure.}
Embed all chunks and queries. Then separate two different questions:

\begin{itemize}[leftmargin=*]
	\item \textbf{Query-local structure, primary generation check}: use only the chunks generated for the current query. This checks whether the synthetic item itself has the intended facet and cluster structure.
	\item \textbf{Full-corpus visibility, primary retrieval check if full-corpus evaluation is reported}: use the full \field{chunks.parquet} corpus. This checks whether the query's gold facets remain visible in the first-stage top-$N$ pool when chunks from other query plans also compete.
\end{itemize}

\paragraph{Query-local checks.}

\begin{itemize}[leftmargin=*]
	\item each required facet has enough gold chunks;
	\item intra-facet similarity is high enough that each facet forms a usable cluster;
	\item inter-facet similarity is lower than intra-facet similarity, so facets do not collapse into one cluster;
	\item the dominant gold cluster is larger or more query-similar than the complementary gold clusters;
	\item hard distractors are close to the query or to a gold facet, but have no gold qrel for the target query.
\end{itemize}

\paragraph{Full-corpus visibility checks.}
If the main evaluation uses full-corpus retrieval:

\begin{itemize}[leftmargin=*]
	\item each required facet has at least one gold chunk in the full-corpus top-$N$ pool;
	\item the top-$N$ pool contains hard distractors or background negatives, so evaluation is not trivial;
	\item chunks from other query plans do not completely swamp the query's gold evidence.
\end{itemize}

\paragraph{Candidate query filters.}

\begin{align}
	\text{FacetPresence}(q)    & = \frac{|\{f: \exists c \in P_N(q), c \in \text{Gold}(q,f)\}|}{|\text{Facets}(q)|}, \\
	\text{TopKDominance}(q,k)  & = \max_f \frac{|\{c \in \text{TopK}(q,k): c \in \text{Cluster}(f)\}|}{k},           \\
	\text{GoldPoolRecall}(q,N) & = \frac{|\text{Gold}(q) \cap P_N(q)|}{|\text{Gold}(q)|}.
\end{align}

Recommended keep rules are structural admission filters, fixed before final testing:

\begin{itemize}[leftmargin=*]
	\item $\text{FacetPresence}(q)=1.0$ for the chosen $N$;
	\item $\text{TopKDominance}(q,10) \geq 0.5$ for at least one common $k$;
	\item $\text{GoldPoolRecall}(q,N) \geq 0.5$ or, better, use pool-relative qrels;
	\item hard distractors are close enough to appear in the candidate pool without receiving positive qrels;
	\item query-local facet clusters remain separated rather than collapsing into one broad cluster.
\end{itemize}

\paragraph{Important warning.}
Do not evaluate against gold chunks that are structurally absent from the candidate pool. Either ensure all gold facets have representatives in the pool or compute recall against \emph{pool-visible gold}. The old MIMIC setup likely suffered from annotation/evaluation mismatch: the annotation pool and full-corpus evaluation pool were not identical.

\stage{Stage 11}{Evaluate Retrieval Methods}

\paragraph{Corpus settings.}
The retrieval corpus is part of the benchmark definition, so different corpus choices should be reported as different settings:

\begin{itemize}[leftmargin=*]
	\item \textbf{Full-corpus retrieval, primary setting}: use the full generated corpus for the split. The corpus contains gold chunks for the current query, its hard distractors, and chunks generated for other query plans.
	\item \textbf{Query-local retrieval, diagnostic ablation}: run the same methods on only the chunks generated for the current query: gold chunks, hard distractors, and query-specific background negatives. This tests reranking behavior after the query-local universe has already been isolated.
\end{itemize}

\paragraph{Retrieval methods.}

\begin{enumerate}[leftmargin=*]
	\item \textbf{Top-$k$}: select the $k$ highest cosine-similarity chunks.
	\item \textbf{MMR}: greedily select chunks balancing query similarity and distance from selected chunks.
	\item \textbf{Facility-location}: greedily optimize a coverage objective over the candidate pool.
\end{enumerate}

\paragraph{Facility-location objective.}
Let $V$ be the candidate pool, $S \subseteq V$ be the selected set, $s_{ij}$ be chunk--chunk similarity, $r_i$ be query relevance, and $w_i=g(r_i)$ be a nonnegative query-relevance weight, for example a clipped or normalized score derived from first-stage similarity. A useful objective is:

\[
	F(S) = \lambda \sum_{i \in V} w_i \max_{j \in S} s_{ij} + (1-\lambda) \sum_{j \in S} r_j.
\]

The weighted coverage term rewards representing high-relevance regions of the candidate pool rather than every dense region equally; the relevance term prevents the method from selecting low-relevance representatives. This matters because the benchmark intentionally contains dense distractor clusters. Tune $\lambda$ and any weighting function parameters on the validation split only.

\paragraph{MMR objective.}

\[
	\operatorname{MMR}(c) = \lambda \operatorname{sim}(q,c) - (1-\lambda)\max_{s \in S}\operatorname{sim}(c,s).
\]

MMR promotes novelty relative to already selected chunks. It can help, but it may select isolated outliers or semantically diverse distractors. The proposed benchmark should make this difference visible: coverage should prefer dense, relevant facet clusters; MMR may overvalue mere distance.

\section{Metrics}

The most important metric is not ordinary document recall. The benchmark is designed around facets.

\begin{longtable}{p{4cm}p{10.9cm}}
	\caption{Recommended evaluation metrics.}\\
	\toprule
	\textbf{Metric} & \textbf{Definition and purpose}\\
	\midrule
	\endfirsthead
	\toprule
	\textbf{Metric} & \textbf{Definition and purpose}\\
	\midrule
	\endhead
	\bottomrule
	\endfoot
	Facet Coverage@k & Fraction of required facets represented by at least one retrieved gold chunk. In this pipeline each required facet is implemented as one gold cluster, so this is also cluster coverage. Primary metric.\\
	Weighted Facet Coverage@k & Average fraction of pool-visible gold chunks retrieved per facet. Captures depth of coverage.\\
	Gold Precision (GP) & Fraction of retrieved chunks that are gold for any facet. Guards against retrieving many irrelevant representatives.\\
	Gold Recall (GR) & Fraction of pool-visible gold chunks retrieved. Report pool-visible recall to avoid penalizing methods for absent gold chunks.\\
	Dominant-Cluster Concentration@k & Fraction of selected chunks drawn from the largest or intentionally dominant gold cluster. Useful for showing the top-$k$ redundancy failure mode.\\
	Redundancy@k & Average number of selected chunks per represented facet or mean pairwise similarity among selected gold chunks. Lower is better if facet coverage is preserved.\\
	Distractor Rate@k & Fraction of selected chunks that are hard distractors. Useful for showing MMR failure modes.\\
	Facet nDCG & Graded ranking metric where each facet contributes gain once or with diminishing returns.\\
	Answer Citation Coverage & Whether a downstream answer generator has enough retrieved citations to support every answer facet.\\
	Facility-location Score & Objective value, reported as a diagnostic, not as the main external metric.\\
	Jaccard vs.\ Top-$k$ & Measures how much a method actually differs from the baseline.\\
\end{longtable}

\paragraph{Primary thesis table.}
The main result should report Facet Coverage@k, Weighted Facet Coverage@k, GP, GR, Dominant-Cluster Concentration@k, and Distractor Rate@k for each method at $k \in \{5,10,20,30\}$.

\paragraph{Expected result pattern.}
The desired pattern is:

\begin{itemize}[leftmargin=*]
	\item facility-location has the highest facet coverage;
	\item top-$k$ has high precision in the dominant facet but lower facet coverage;
	\item MMR has higher diversity than top-$k$ but worse GP and higher distractor rate than facility-location;
	\item the gap is largest at small or moderate $k$, e.g.\ $k=5$ and $k=10$;
	\item at high $k$, top-$k$ may catch up because it eventually retrieves many clusters by volume.
\end{itemize}

\section{Concrete Generation Example}

\subsection{Hidden Query Plan}

\begin{lstlisting}[style=promptstyle]
Primary condition:
  encephalitis or myelitis

Required comparison groups:
  A: uncomplicated diabetes
  B: age over 75

Required clinical axes:
  1. treatment duration
  2. rehabilitative outcome

Gold facets:
  F1: diabetes + treatment duration
  F2: diabetes + rehabilitative outcome
  F3: elderly + treatment duration
  F4: elderly + rehabilitative outcome
\end{lstlisting}

\subsection{Generated Natural Query}

\begin{quote}
	For patients diagnosed with encephalitis or myelitis, how do treatment durations and rehabilitative outcomes differ between those with uncomplicated diabetes and elderly patients over age 75?
\end{quote}

\subsection{Gold Cluster Plan}

\begin{longtable}{p{2.6cm}p{4.2cm}p{4.1cm}p{3.5cm}}
	\caption{Example gold and distractor clusters for one query.}\\
	\toprule
	\textbf{Cluster} & \textbf{Contents} & \textbf{Gold role} & \textbf{Target size}\\
	\midrule
	\endfirsthead
	\toprule
	\textbf{Cluster} & \textbf{Contents} & \textbf{Gold role} & \textbf{Target size}\\
	\midrule
	\endhead
	\bottomrule
	\endfoot
	C1 & encephalitis + uncomplicated diabetes + antiviral duration & dominant gold, facet F1 & 18 chunks\\
	C2 & encephalitis + uncomplicated diabetes + rehab outcome & complementary gold, facet F2 & 8 chunks\\
	C3 & encephalitis + age $>75$ + treatment duration & complementary gold, facet F3 & 8 chunks\\
	C4 & encephalitis + age $>75$ + rehab outcome & complementary gold, facet F4 & 8 chunks\\
	D1 & encephalitis + no diabetes + age 40--64 & hard distractor & 20 chunks\\
	D2 & elderly + stroke + rehab outcome & hard distractor & 20 chunks\\
	D3 & diabetes + pneumonia + treatment duration & hard distractor & 20 chunks\\
	D4 & general encephalitis background & generic distractor & 5 chunks\\
\end{longtable}

\subsection{Expected Retrieval Behavior}

Top-$k$ should retrieve many C1 chunks because the natural query contains condition, diabetes, and treatment-duration terms. MMR should retrieve from C1 but then jump toward semantically distinct clusters, including D2 or D3. Facility-location should select representatives from C1--C4 because they are dense regions in the candidate pool and each contributes coverage.

\section{Recommended Ablations}

\begin{enumerate}[leftmargin=*]
	\item \textbf{No distractors}: verifies that all methods improve when the task is easy.
	\item \textbf{No redundancy}: removes the condition under which facility-location should help; expected to shrink the gap.
	\item \textbf{No structured fact grounding}: use naive LLM-generated notes and show quality/gold-label degradation.
	\item \textbf{MIMIC-seeded vs.\ fully synthetic}: compare style realism and retrieval geometry.
	\item \textbf{Different embedding models}: general embedding vs.\ biomedical embedding.
	\item \textbf{Single-query vs.\ subquery fusion}: tests whether explicit multi-aspect query decomposition competes with reranking.
	\item \textbf{Facet count}: 2, 3, 4, and 5 facets per query.
	\item \textbf{Distractor strength}: easy, medium, hard, and outlier distractors.
\end{enumerate}

\section{Prompt and Model Strategy}

\paragraph{Generator model.}
Do not assume that the best answering model is the best generator. AGORABench shows that model generation ability and task-solving ability can diverge \citep{kim2025agorabench}. If resources permit, compare two generator models on a small pilot batch before freezing the generation pipeline, then choose the one with better structured validity, diversity, and clinical plausibility.

\paragraph{Temperatures.}
Use low temperature for structured facts and validation. Use moderate temperature for narrative variation. DATAGEN reports diversity benefits from generation settings and attribute guidance \citep{huang2025datagen}; however, benchmark labels should remain deterministic.

\paragraph{Prompt optimization.}
If generation quality is unstable, an optional later phase can optimize prompts for two objectives: clinical fidelity and diversity. This is inspired by multi-objective prompt learning \citep{pastrian2026evolutionary}, but it is not necessary for the first implementation.

\section{Implementation Steps}

\begin{enumerate}[leftmargin=*]
	\item Create \field{ontology.yaml} with conditions, concepts, comorbidities, demographics, interventions, outcomes, and allowed combinations.
	\item Create \field{profile\_templates.yaml} or an equivalent latent-profile table for schema-first generation and split control.
	\item Generate \field{query\_plans.parquet}: sample query family + condition + 3--5 objective facets + difficulty label.
	\item Generate \field{clinical\_facts.parquet}: structured atomic facts for each facet and distractor.
	\item Validate facts using schema, clinical plausibility, concept-consistency, and subgroup-balance checks.
	\item Generate \field{chunks.parquet}: narrative text from facts.
	\item Apply role-aware deduplication, non-copying checks, and group-balance enforcement.
	\item Generate \field{queries.parquet}: natural query text plus hidden subqueries and optional surface variants.
	\item Generate \field{qrels.parquet}: directly from \field{chunk\_id}--\field{fact\_id}--\field{facet\_id} mappings.
	\item Generate \field{gold\_answers.parquet}: one answer per query with facet citations.
	\item Write \field{generation\_audit.parquet}: prompt hashes, model versions, decoding parameters, seeds, validators, rejection reasons.
	\item Freeze validation/test splits by query family/profile before final threshold tuning.
	\item Embed chunks and queries.
	\item Compute candidate pools and geometry diagnostics.
	\item Filter queries based on predeclared facet presence, top-$k$ dominance, pool-visible gold, distractor proximity, and query-local cluster separation.
	\item Tune $\lambda$ on validation.
	\item Report final metrics on test.
\end{enumerate}

\section{Quality Report to Include in the Thesis}

The benchmark should include a transparent generation report. Following the concerns raised in biomedical synthetic data surveys \citep{rao2026scoping,montenegro2025medicalsynthetic}, report:

\begin{itemize}[leftmargin=*]
	\item number of generated facts, chunks, queries, facets, and distractors;
	\item rejection rates by validation stage;
	\item average chunks per facet;
	\item distribution of query types and difficulty levels;
	\item average within-facet and across-facet embedding similarity;
	\item no-context answerability rate;
	\item prompt hashes, model identifiers, decoding settings, schema versions, and random seeds;
	\item profile-level and subgroup-level split distributions;
	\item non-copying/authenticity check results when real examples seed the generator;
	\item cluster-quality diagnostics such as intra-facet similarity, inter-facet separation, and silhouette scores;
	\item human audit size and agreement if available;
	\item examples of accepted and rejected generations;
	\item limitations of synthetic clinical realism.
\end{itemize}

\section{Threats to Validity}

\paragraph{Synthetic bias.}
The corpus is generated to favor coverage. This is acceptable if the thesis frames the benchmark as a controlled stress test. It would be problematic only if presented as evidence that facility-location always improves real clinical RAG.

\paragraph{Authenticity and privacy.}
If MIMIC-IV snippets or guideline text are used as style seeds, generated chunks may accidentally copy source phrasing. Mitigate with abstraction prompts, overlap checks, source hashes, and manual spot audits. The final corpus should contain synthetic evidence, not republished patient text.

\paragraph{Medical realism.}
LLM-generated facts may be clinically imperfect. Mitigate with constrained templates, concept grounding, decision-pathway grounding, and validation. Avoid using the dataset for clinical recommendations.

\paragraph{Embedding dependence.}
The result may depend on the embedding model. Report experiments with at least one general-purpose and one biomedical embedding model if feasible.

\paragraph{Query filtering bias.}
Filtering for geometric properties can look like cherry-picking. Mitigate by defining filters before final evaluation, applying them uniformly, and reporting how many queries were rejected for each reason.

\paragraph{Gold-label leakage.}
If narrative chunks accidentally mention all facets, the task becomes single-source. Reject chunks that support too many facets unless intentionally creating an easy split.

\paragraph{MMR parameter fairness.}
MMR must be tuned on the same validation split as facility-location. Otherwise, the comparison is not fair.

\section{Final Recommended Experimental Claim}

The strongest thesis claim is:

\begin{quote}
	We construct a controlled synthetic clinical RAG benchmark in which each query is generated from an explicit multi-facet evidence plan. The corpus contains redundant relevant clusters, complementary gold clusters, and hard distractors. Under this controlled multi-aspect setting, facility-location reranking improves facet-level evidence coverage over similarity top-$k$ and reduces the distractor sensitivity of dispersion-only MMR.
\end{quote}

This claim is narrower than ``facility-location is better for medical RAG'', but it is much more defensible and experimentally clean.

\section{Minimal First Version}

If time is limited, implement the smallest useful version:

\begin{itemize}[leftmargin=*]
	\item 4 primary conditions;
	\item 100 queries;
	\item 4 facets per query;
	\item 10 gold chunks per facet;
	\item 30 distractors per query;
	\item 10,000--20,000 total chunks;
	\item top-$k$, MMR, facility-location at $k \in \{5,10,20\}$;
	\item Facet Coverage@k, Weighted Facet Coverage@k, GP, GR, Dominant-Cluster Concentration@k, and Distractor Rate@k.
\end{itemize}

After this works, scale to more conditions and query types.

\bibliographystyle{plainnat}
\bibliography{synthetic-rag-benchmark-pipeline}

\end{document}
